% Generated draft: requires mathematical and source review.
\documentclass[11pt]{article}
\usepackage[margin=25mm]{geometry}
\usepackage{fontspec}
\setmainfont{Latin Modern Roman}
\usepackage{amsmath,amssymb,mathtools,unicode-math}
\setmathfont{Latin Modern Math}
\usepackage{enumitem,needspace,xurl,hyperref}
\setlength{\parindent}{0pt}
\setlength{\parskip}{5pt}
\setlength{\emergencystretch}{3em}
\newcommand{\N}{\mathbb{N}}
\newcommand{\Z}{\mathbb{Z}}
\newcommand{\Q}{\mathbb{Q}}
\newcommand{\R}{\mathbb{R}}
\newcommand{\C}{\mathbb{C}}
\newcommand{\F}{\mathbb{F}}
\newcommand{\A}{\mathbb{A}}
\newcommand{\PP}{\mathbb{P}}
\newcommand{\E}{\mathbb{E}}
\newcommand{\eps}{\varepsilon}
\newcommand{\dd}{\,\mathrm{d}}
\begin{document}
\section*{UnsolvedMath ID: 8800041}
\textbf{Source problem code:} catalogue entry
\subsection{Definitions and mathematical statement}
Let $p$ be prime and let $E/\mathbb F_p$ be an elliptic curve, with finite group order $N_E=\#E(\mathbb F_p)$. For $y\geq2$, an integer $N\geq1$ is $y$-smooth if every prime divisor of $N$ is at most $y$. Equivalently, $P^+(N)\leq y$, where $P^+(N)$ is its largest prime factor and $P^+(1)=1$.

Prove a sufficiently strong lower bound for the number of curves with $P^+(N_E)\leq y$ in the intended smoothness regime. For a specified probability law $\nu_p$ on elliptic-curve isomorphism classes, a proportion is
\[
S(p,y)=\Pr_{E\sim\nu_p}\bigl[P^+(\#E(\mathbb F_p))\leq y\bigr].
\]
One may instead count coefficient models or use automorphism-weighted class counts; the source does not choose among them.

The words ``smooth'' and ``sufficiently many'' leave the essential quantitative problem unspecified: the bound $y=y(p)$, the desired lower bound on $S(p,y)$ or a raw count, and uniformity over primes are absent. A statement for every sufficiently large prime, one for almost all primes, and an average-over-primes statement are different targets. No unrestricted positive-density claim is invented. The supplied summary reports nonoptimal smoothness estimates from elliptic-curve/Hasse-interval and smooth-number methods, but gives no exact theorem or benchmark.
\subsection{Short English statement}
Show that many elliptic curves over a prime field have group orders with only small prime factors. Fix the smoothness cutoff, curve-counting measure, and quantitative lower bound and prime-uniformity requirements.
\subsection{Sources}
\begin{itemize}
\item Igor Shparlinski, IPAM \emph{Workshop on Number Theory and Cryptography --- Open Problems} (2006), attachment 6224, PDF page 21, schedule source \url{https://www.ipam.ucla.edu/programs/workshops/workshop-i-number-theory-and-cryptography-open-problems/?tab=schedule}; accessed 2026-07-29. The supplied provenance records a PDF/legacy-PPT visual audit.
\item Related attributions in the supplied summary: Lenstra, Hasse-interval distribution methods; Luca--Shparlinski, smooth-order counting; Gaudry, Hess, and Smart, related Weil-descent/index-calculus smoothness. No complete publication identifiers or precise hypotheses are supplied.
\item Local registry \texttt{data/unsolvedmath/problems.json}, permanent UnsolvedMath ID $8800041$, source code catalogue entry, \emph{Elliptic curves with smooth group order}. Dataset attribution: Ulam AI and the UnsolvedMath Contributors.
\end{itemize}
\end{document}
